\documentclass[12pt,letterpaper]{article}



%%
%% Required packages:
%%
\usepackage{etex}
\usepackage{amsmath}
\usepackage{amssymb}
\usepackage{amstext}
\usepackage{amsthm}
\usepackage{fancyhdr,fancybox}
\usepackage{mathrsfs} % need for \mathscr command
\usepackage{multicol}
\usepackage[normalem]{ulem}
%\usepackage{pictex,pictexplus}
% \usepackage{pictexwd}
\usepackage{rotating}
\usepackage{url}
\usepackage{xfrac}
\usepackage{booktabs}
\usepackage{color,soul}
\usepackage{comment}

\usepackage{hyperref}
\hypersetup{
    colorlinks=true,     % false: boxed links; true: colored links
    linkcolor=blue,      % color of internal links
    citecolor=blue,      % color of links to bibliography
    filecolor=blue,      % color of file links
    urlcolor=blue        % color of external links
}


\usepackage{tikz}
\usetikzlibrary{backgrounds,calc}

%%
%% Common formatting commands
%%
\setlength{\parindent}{0in}
\setlength{\textwidth}{7in}
\setlength{\evensidemargin}{-0.25in}
\setlength{\oddsidemargin}{-0.25in}
\setlength{\parskip}{.5\baselineskip}
\setlength{\topmargin}{-0.5in}
\setlength{\textheight}{9in}

%               Problem and Part
\newcounter{problemnumber}
\newcounter{partnumber}
\newcounter{subpartnumber}
\newcommand{\Problem}{%
  \refstepcounter{problemnumber}%
  \setcounter{partnumber}{0}%
  \setcounter{subpartnumber}{0}%
  \item[\makebox{\hfil\textbf{\theproblemnumber.}\hfil}]%
  }
\newcommand{\InProblem}[2][1.6]{%
  \refstepcounter{problemnumber}%
  \setcounter{partnumber}{0}%
  \setcounter{subpartnumber}{0}%
  \textbf{\theproblemnumber.}\ \ \parbox[t]{#1in}{#2}%
}
\newcommand{\InSmallProblem}[1]{\InProblem[1.05]{#1}}
\newcommand{\NoProblem}[1]{%
  \mbox{\hfil\textbf{#1}\hfil}%
  }
\newcommand{\UnProblem}{%
  \refstepcounter{problemnumber}%
  \setcounter{partnumber}{0}%
  \setcounter{subpartnumber}{0}%
  \mbox{\hfil\textbf{\theproblemnumber.}\hfil}%
  }
\newcommand{\InFlexProblem}[1]{%
  \refstepcounter{problemnumber}%
  \setcounter{partnumber}{0}%
  \setcounter{subpartnumber}{0}%
  \textbf{\theproblemnumber.}\ \ #1%
}
%%  Part:
\newcommand{\Part}{%
  \renewcommand{\thepartnumber}{(\alph{partnumber})}%
  \refstepcounter{partnumber}
  \setcounter{subpartnumber}{0}%
  \item[(\alph{partnumber})]%
}
\newcommand{\InPart}[2][1.75]{%
  \renewcommand{\thepartnumber}{(\alph{partnumber})}%
  \refstepcounter{partnumber}%
  \setcounter{subpartnumber}{0}%
  (\alph{partnumber})\ \ \parbox[t]{#1in}{#2}%
}
\newcommand{\UnPart}{%
  \refstepcounter{partnumber}%
  \renewcommand{\thepartnumber}{(\alph{partnumber})}%
  \setcounter{subpartnumber}{0}%
  (\alph{partnumber})%
}
\newcommand{\InLargePart}[1]{\InPart[3.05]{#1}}
\newcommand{\InFullPart}[1]{\InPart[6.2]{#1}}
\newcommand{\InSmallPart}[1]{\InPart[1.05]{#1}}
%% SubPart:
\newcommand{\SubPart}{%
  \refstepcounter{subpartnumber}%
  \item[(\roman{subpartnumber})]%
  }
\newcommand{\InSubPart}[2][2.25]{%
  \stepcounter{subpartnumber}%
  (\roman{subpartnumber})\ \ \parbox[t]{#1in}{#2}%
}
\newcommand{\InSmallSubPart}[1]{\InSubPart[1.5]{#1}}
\newcommand{\InTinySubPart}[1]{%
  \stepcounter{subpartnumber}%
  (\roman{subpartnumber})\ \ \makebox{#1}%
}
\newcommand{\InMySubPart}[2][2.25]{%
  \stepcounter{subpartnumber}%
  \parbox[t]{#1in}{\makebox[0.3in][l]{(\roman{subpartnumber})}\ \ {#2}}%
}
\newcommand{\TrueFalse}{\textbf{T}\ \ \textbf{F}\ \ }
\newcommand{\TFPart}[1]{% 
  \stepcounter{partnumber}%
  \setcounter{subpartnumber}{0}%
  \item[(\alph{partnumber})] \TrueFalse\parbox[t]{5.5in}{#1}}

\newcommand{\Points}[1]{(#1 points) \ }
\newcommand{\Point}[1]{(#1 point) \ }


%% For Answers:
\newcounter{answernumber}
\newcommand{\Answer}{%
  \refstepcounter{answernumber}%
  \setcounter{partnumber}{0}%
  \setcounter{subpartnumber}{0}%
  \item[\makebox{\hfil\textbf{\theanswernumber.}\hfil}]%
  }


% Setting up the headers for the first page
%\chead{} % will default to what is typed in the file.
\fancyhead[L]{\textsc{Math 22a}}
\fancyhead[R]{\textsc{Fall 2024}}
\fancyfoot[R]{
	\vspace*{\fill}\footnotesize\textcopyright{} \emph{Harvard Math 22a}	
}
\renewcommand{\headrulewidth}{0pt}

%\fancyhead[C]{}
%	\chead{}	
%	\lhead{}
%	\rhead{}
%	\cfoot{}


% Putting copyright on the footer of each page
%\fancypagestyle{secondpage}{
%	\fancyhead[C]{TESTing again} % change this to be blank
%		%\chead{}	
%	\fancyhead[L]{bears} % change this to be blank
%	\fancyhead[R]{dogs} % change this to be blank
%	\fancyfoot[R]{ %keeps the footer the same
%		\vspace*{\fill}\footnotesize\textcopyright{} \emph{Harvard Math 22a}	
%	}
%}
%\renewcommand{\headrulewidth}{0pt}
%\pagestyle{fancy}
\pagestyle{empty}


%\lhead{}
%\rhead{}
%\rfoot{\vspace*{\fill}\footnotesize\textcopyright{} \emph{Harvard Math 22a}}
%\renewcommand{\headrulewidth}{0pt}
%\pagestyle{fancy}




%%
%% Common shortcut commands:
%%
\newcommand{\ds}{\displaystyle}
\newcommand{\sgn}{\operatorname{sgn}}
\renewcommand{\Pr}{\operatorname{P}}
\newcommand{\CPr}[3][{}]{\Pr_{\text{#1}}\left(#2\,|\,#3\right)}

\newcommand{\C}{\mathbb{C}}
\newcommand{\R}{\mathbb{R}}
\newcommand{\Z}{\mathbb{Z}}
\newcommand{\N}{\mathbb{N}}
\newcommand{\Q}{\mathbb{Q}}
\newcommand{\D}{\mathbb{D}}

\newcommand{\rref}{\operatorname{rref}}
\newcommand{\im}{\operatorname{im}}
\newcommand{\rank}{\operatorname{rank}}
\newcommand{\diag}{\operatorname{diag}}
\newcommand{\Span}{\operatorname{span}}
%\renewcommand{\vec}[1]{\mathbf{#1}}
\newcommand{\charpoly}{\mathscr{P}}
\newcommand{\nicefrac}[2]{\sfrac{#1}{#2}} % using xfrac package
\newcommand{\topology}{\mathcal{T}}
\newcommand{\Inn}{\operatorname{Inn}}

\newcommand{\blankbox}[2]{\fbox{\rule{#1}{0in}\rule{0in}{#2}}}

\newenvironment{myindentpar}[1]%
 {\begin{list}{}%
         {\setlength{\leftmargin}{#1}
          \setlength{\rightmargin}{\leftmargin}}%
         \item[]%
 }
 {\end{list}}
\newtheorem{theorem}{Theorem}
\newtheorem*{NStheorem}{Theorem}
\newtheorem{cor}[theorem]{Corollary}
\newtheorem*{claim}{Claim}
\newtheorem{defn}{Definition}

\newenvironment{amatrix}[1]{%
  \left(\begin{array}{@{}*{#1}{c}|c@{}}
}{%
  \end{array}\right)
}

%--------grstep
% For denoting a Gauss' reduction step.
% Use as: \grstep{\rho_1+\rho_3} or \grstep[2\rho_5 \\ 3\rho_6]{\rho_1+\rho_3}
\newcommand{\grstep}[2][\relax]{%
   \ensuremath{\mathrel{
       {\mathop{\longrightarrow}\limits^{#2\mathstrut}_{
                                     \begin{subarray}{l} #1 \end{subarray}}}}}}
\newcommand{\swap}{\leftrightarrow}


\usetikzlibrary{arrows}
\usepackage{wrapfig}

\makeatletter
\renewcommand*\env@matrix[1][*\c@MaxMatrixCols c]{%
	\hskip -\arraycolsep
	\let\@ifnextchar\new@ifnextchar
	\array{#1}}
\makeatother

\def\env@matrix{\hskip -\arraycolsep
	\let\@ifnextchar\new@ifnextchar
	\array{*\c@MaxMatrixCols c}}

\chead{\Large\textbf{Complex Eigenvalues, $\vec\S$5.5}}% 

\begin{document}
\thispagestyle{fancy}

Recall from previously:

\begin{itemize}
\item If $A$ is an $n \times n$ matrix, then $\lambda \in \mathbb{R}$ is an {\bf eigenvalue} of $A$ if there exists a non-zero vector $v$, called an {\bf eigenvector} corresponding to $\lambda$, such that $$Av=\lambda v.$$
\item If $\lambda$ is an eigenvalue of an $n \times n$ matrix $A$, then the set of all solutions $v$ to $Av = \lambda v$ is the {\bf eigenspace} of $A$ corresponding to $\lambda$.
\item {\bf Theorem 1:} The eigenvalues of a triangular matrix are the diagonal entries.
\item {\bf Theorem 2:} If $v_1, \dots, v_p$ are eigenvectors corresponding to distinct eigenvalues $\lambda_1, \dots, \lambda_p$ of an $n\times n$ matrix $A$, then $v_1, \dots, v_p$ are linearly independent.
\item If $A$ is an $n\times n$ matrix, then the characteristic polynomial of $A$, denoted by $p_A(\lambda)$, is given by $$p_A(\lambda)=\text{det}(A-\lambda I_n).$$
\item $\lambda$ is an eigenvalue of an $n\times n$ matrix $A$ if and only if $\lambda$ is a zero of the characteristic polynomial.
\item If $A$ and $B$ are similar matrices, then they have the same characteristic polynomials and hence the same eigenvalues (with the same algebraic multiplicities). 
\item An $n\times n$ matrix $A$ is {\bf diagonalizable} if $A$ is similar to a diagonal matrix.
\item {\bf Diagonalization Theorem:} An $n \times n$ matrix $A$ is diagonalizable if and only if $A$ has $n$ linearly independent eigenvectors. Moreover, $A=PDP^{-1}$ with $D$ diagonal if and only if the columns of $P$ are $n$ linearly independent eigenvectors of $A$. In this case the diagonal entires of $D$ are the corresponding eigenvalues of $A$.
\item {\bf Theorem 6:} An $n \times n$ matrix $A$ with $n$ distinct eigenvalues is diagonalizable.
\item 
{\bf Theorem 7:} Let $A$ be an $n\times n$ matrix with distinct eigenvalues $\lambda_1, \dots, \lambda_p$.
\begin{enumerate}
\item For $1\leq k \leq p$, the dimension of the eigenspace of $\lambda_k$ is at most the algebraic multiplicity of the eigenvalue $\lambda_k$.
\item The matrix $A$ is diagonalizable if and only if the sum of the dimensions of the eigenspaces is $n$ if and only if the characteristic polynomial factors into linear factors and the dimension for each eigenspace is exactly the algebraic multiplicity.
\item If $A$ is diagonalizable and $\mathcal{B}_k$ is a basis for the eigenspace corresponding to $\lambda_k$, then $\mathcal{B}_1, \dots, \mathcal{B}_p$ is an eigenbasis for $\mathbb{R}^n$.
\end{enumerate}
\end{itemize}

\newpage


\begin{enumerate}

\Problem Let $A=\begin{bmatrix} 0 & -1\\ 1 & 0\\ \end{bmatrix}$. Find eigenvalues and eigenspaces.

\vfill

\Problem Let $A$ be an $n \times n$ matrix with real entries. Suppose $\lambda$ is an eigenvalue of $A$, then is $\overline{\lambda}$ an eigenvalue of $A$.

\vfill

\newpage

\Problem Let $C=\begin{bmatrix} a & -b \\ b & a\\ \end{bmatrix}$ where $a, b$ are real and not both zero, then find the eigenvalues of $C$.

\vfill

{\bf Theorem:} Let $A$ be a real 2 by 2 matrix with a complex eigenvalue $\lambda=a-bi$ ($b\neq0$) and an associated eigenvector $v$ in $\mathbb{C}^2$. Then $A=PCP^{-1}$ where $P=[ \text{Re}(v) \; \text{Im}(v) ]$ and $C=\begin{bmatrix} a & -b\\ b & a\\ \end{bmatrix}$.

\Problem Prove the theorem.

\vfill

\newpage

\Problem Let $A=\begin{bmatrix} 3 & -2\\ 4 & -1\\ \end{bmatrix}$, then find $P$ and $C$ such that $A=PCP^{-1}$ as in the previous theorem.

%\Problem Consider $x'(t) = \begin{bmatrix} 1 & 2\\ 3 & 2\\ \end{bmatrix} x(t)$. Solve the differential equation.

%\vfill

%\Problem Consider $x'(t) = \begin{bmatrix} 0 & -1\\ 1 & 0\\ \end{bmatrix} x(t)$. Solve the differential equation.

%\vfill

\end{enumerate}

\begin{comment}
\newpage
\chead{\Large\textbf{Complex Eigenvalues-- Answers / Solutions}}%
\lhead{}
\rhead{}
\thispagestyle{fancy}

\begin{enumerate}
\Answer We find the eigenvalues using the characteristic polynomial; namely, $$p_A(\lambda)=\det \begin{bmatrix} -\lambda & -1\\ 1 & -\lambda\\ \end{bmatrix}= \lambda^2+1.$$ Thus the zeros of the characteristic polynomial are $\lambda=\pm i.$ We find the eigenspaces now.

\underline{$\lambda=i$.} 

Here we get $$A-iI_2 = \begin{bmatrix} -i & -1\\ 1 & -i\\ \end{bmatrix},$$ and hence $E_{i}=N(A-iI_2)= \text{Span}\left \{\begin{bmatrix}1\\ -i\\ \end{bmatrix} \right\}.$

\underline{$\lambda=-i$.} 

Here we get $$A+iI_2 = \begin{bmatrix} i & -1\\ 1 & i\\ \end{bmatrix},$$ and hence $E_{-i}=N(A+iI_2)= \text{Span}\left \{\begin{bmatrix}1\\ i\\ \end{bmatrix} \right\}.$

Thus we can digaonalize $A$ over the complex numbers as $$A=PDP^{-1}=\begin{bmatrix} 1 & 1\\ -i & i\\ \end{bmatrix} \begin{bmatrix} i & 0\\ 0 & -i\\ \end{bmatrix}\begin{bmatrix} 1 & 1\\ -i & i\\ \end{bmatrix}^{-1}.$$

\Answer Let $A$ be an $n\times n$ matrix with real entries, then the complex conjugate of $A$ is still just $A$;i.e. $\overline{A}=A$. Let $\lambda$ be an eigenvalue of $A$ with eigenvector $v$. Then $$\overline{\lambda} \overline{v} = \overline{\lambda v} = \overline{A v}=\overline{A} \overline{v}= A \overline{v},$$ and hence $\overline{\lambda}$ is also an eigenvalue of $A$ with eigenvector $\overline{v}$. 
\Answer We again use the characteristic polynomial to find the eigenvalues of $C$; namely, $$p_{C}(\lambda)= \det{(C- \lambda I_2)}=\det\begin{bmatrix}a-\lambda & -b\\ b &a-\lambda \end{bmatrix}=a^2+b^2-2a\lambda +\lambda^2.$$ Thus using the quadratic formula, we get the zeros of the characteristic polynomial are $\lambda=a\pm bi$.

Expressing $\lambda=a+bi$ in polar form, we get $\lambda= r (\cos{\theta}+i\sin{\theta}),$ and hence $a=r\cos{\theta}$ and $b= r\sin{\theta}$. Thus $$C=\begin{bmatrix} a & -b\\ b & a\\ \end{bmatrix}= \begin{bmatrix} r \cos \theta & -r \sin \theta\\ r \sin \theta & r \cos \theta\\ \end{bmatrix}= \begin{bmatrix} r & 0 \\ 0 & r\\ \end{bmatrix} \begin{bmatrix} \cos \theta & - \sin \theta\\ \sin \theta & \cos \theta\\ \end{bmatrix}.$$ Thus we can regard multiplying by matrices of the form $C$ as rotating by $\theta$ and scaling by $r$!

\Answer

\begin{proof}
	Let $A$ be a $2\times 2$ matrix with real entries. Suppose $A$ has eigenvalue $\lambda= a-bi$ with $b\neq 0$ and associated complex eigenvector $v$. Define the change of coordinates matrix $P$ by $P=[ \text{Re}{(v)} \; \text{Im}{(v)}]$. We want to prove $A= PCP^{-1}$ where $C$ has the desired form. We will prove that $AP=PC$. We begin by finding $A\text{Re}{(v)}$ and $A\text{Im}{(v)}$. Now 
	\begin{align*}
	A \text{Re}(v)+iA\text{Im}(v)&=A(\text{Re}(v)+\text{Im}(v))\\
	&=Av=\lambda v=(a-bi)(\text{Re}(v)+i\text{Im}(v))\\
	&=(a\text{Re}(v)+b\text{Im}(v))+i(-b\text{Re}(v)+a \text{Im}(v))\\
	&= P\begin{bmatrix} a\\ b\\ \end{bmatrix} +i P \begin{bmatrix} -b\\ a\\ \end{bmatrix},
\end{align*}
and hence $A\text{Re}(v)= P\begin{bmatrix} a\\ b\\ \end{bmatrix}$ and $A\text{Im}(v)= P\begin{bmatrix} -b\\ a\\ \end{bmatrix}$. Thus $$AP= [A\text{Re}(v) \; \; A \text{Im}(v)] =\left[ P\begin{bmatrix} a\\ b\\ \end{bmatrix}\;\; P\begin{bmatrix} -b\\ a\\ \end{bmatrix}\right]=P\begin{bmatrix} a & -b\\ b & a\\ \end{bmatrix}=PC.$$
\end{proof}
\Answer Finding the eigenvalues and eigenspaces as usual, we find the eigenvalue are $1\pm 2i$ and the eigenspaces are $$E_{1+2i}=\text{Span}\left\{\begin{bmatrix}1+i\\ 2\\ \end{bmatrix} \right\},$$ and $$E_{1-2i}=\text{Span}\left\{\begin{bmatrix}1-i\\ 2\\ \end{bmatrix} \right\}.$$ Using the previous theorem, we take $\lambda=1-2i$ and hence $a=1$ and $b=2$ in the Theorem. Take $P=[\text{Re}(v) \; \text{Im}(v)] = \begin{bmatrix} 1 & -1\\ 2 & 0\\ \end{bmatrix}$. Using the previous Theorem, we get $$A=PCP^{-1}= \begin{bmatrix} 1 & -1\\ 2 & 0\\ \end{bmatrix} \begin{bmatrix} 1 & -2\\ 2 & 1\\ \end{bmatrix} \begin{bmatrix} 1 & -1\\ 2 & 0\\ \end{bmatrix}^{-1}.$$

\end{enumerate}  

\end{comment}
\end{document}


%%% Local ables: 
%%% mode: latex
%%% TeX-master: t
%%% End: 
